\section{Space Complexity --- Machines --- ARMProof}
%%% generated by the blueprint pipeline from TCSlib.Complexity.SpaceComplexity.Machines.ARMProof
\subsection*{Space Complexity --- Machines --- ARMProof: Introduction}

Declarations of the Lean module \texttt{TCSlib.Complexity.SpaceComplexity.Machines.ARMProof}.

\subsection*{Space Complexity --- Machines --- ARMProof: Results}

\begin{definition}[Reachability under an invariant]
\lean{Complexity.LogProg.AReach}
\label{Complexity.LogProg.AReach}
\leanok
\uses{Complexity.LogProg.AConf, Complexity.LogProg.ARM, Complexity.LogProg.arun}
Let $A$ be a program, let an oracle assign a Boolean answer to each decider and bit
string, let $x$ be an input, and let $G$ be a predicate on configurations. A
configuration $a'$ is reachable from $a$ in some finite number $T$ of steps if the run
from $a$ on $x$ reaches $a'$ at time $T$ and every configuration at a time $t<T$
satisfies $G$.
\end{definition}

\begin{definition}[Halting with a specified answer]
\lean{Complexity.LogProg.AHalt}
\label{Complexity.LogProg.AHalt}
\leanok
\uses{Complexity.LogProg.AConf, Complexity.LogProg.ARM, Complexity.LogProg.arun}
Let $A$ be a program, let an oracle assign a Boolean answer to each decider and bit
string, let $x$ be an input, and let $G$ be a predicate on configurations. The run from
configuration $a$ halts with answer $b$ if, for some natural-number time $T$, its state
component is empty and its output component is $b$ at time $T$, and every configuration
at a time $t<T$ satisfies $G$.
\end{definition}

\begin{lemma}[Additivity of abstract run lengths]
\lean{Complexity.LogProg.arun_add}
\label{Complexity.LogProg.arun_add}
\leanok
\uses{Complexity.LogProg.AConf, Complexity.LogProg.arun}
\difficulty{3}
For any register program $A$, oracle function $o$, input $x$, initial configuration $a$,
and natural numbers $s,t$, running $A$ for $s+t$ steps from $a$ gives the same
configuration as running it first for $s$ steps and then for $t$ more steps.
\end{lemma}

\begin{lemma}[Reflexivity of guarded reachability]
\lean{Complexity.LogProg.AReach.refl}
\label{Complexity.LogProg.AReach.refl}
\leanok
\uses{Complexity.LogProg.AConf, Complexity.LogProg.AReach}
\difficulty{2}
For any natural numbers $m,d$, type $\Lambda$, register program $A : \Lambda \to
\mathrm{Ins}(m,d,\Lambda)$, oracle $O : \mathrm{Fin}(d) \to \{0,1\}^* \to \{0,1\}$,
input bit string $x$, predicate $G$ on configurations, and configuration $a$, there is a
time $T$ such that the run from $a$ is at $a$ after $T$ steps and $G$ holds at every
configuration at each time $t<T$.
\end{lemma}

\begin{lemma}[Transitivity of Constrained Reachability]
\lean{Complexity.LogProg.AReach.trans}
\label{Complexity.LogProg.AReach.trans}
\leanok
\uses{Complexity.LogProg.AConf, Complexity.LogProg.AReach, Complexity.LogProg.arun_add}
\difficulty{5}
Fix a program, oracle, input, and predicate $G$ on configurations. If there is a finite
run from $a$ to $b$ and a finite run from $b$ to $c$, with every configuration before
each run's endpoint satisfying $G$, then there is a finite run from $a$ to $c$ with
every configuration before its endpoint satisfying $G$.
\end{lemma}

\begin{lemma}[Concatenation of guarded reachability and halting]
\lean{Complexity.LogProg.AReach.halt}
\label{Complexity.LogProg.AReach.halt}
\leanok
\uses{Complexity.LogProg.AConf, Complexity.LogProg.AHalt, Complexity.LogProg.AReach, Complexity.LogProg.arun_add}
\difficulty{5}
Let $A$ be a register program with oracle, let $x$ be an input, and let $G$ be a
predicate on configurations. If a run from configuration $a$ reaches configuration $b$
with $G$ holding at every configuration before reaching $b$, and a run from $b$ halts
with answer $r$ while satisfying $G$ at every configuration before halting, then the run
from $a$ halts with answer $r$ and satisfies $G$ at every configuration before halting.
\end{lemma}

\begin{lemma}[One-step guarded reachability]
\lean{Complexity.LogProg.AReach.single}
\label{Complexity.LogProg.AReach.single}
\leanok
\uses{Complexity.LogProg.AConf, Complexity.LogProg.AReach, Complexity.LogProg.astep}
\difficulty{4}
Let $A$ be a register program with oracle, let $x$ be an input, and let $G$ be a
predicate on configurations. If $G$ holds of configuration $a$, then the one-step
successor of $a$ is reachable from $a$ in one step, with $G$ holding at every
configuration before the endpoint.
\end{lemma}

\begin{lemma}[Prepending a Reachability Step]
\lean{Complexity.LogProg.AReach.step}
\label{Complexity.LogProg.AReach.step}
\leanok
\uses{Complexity.LogProg.AConf, Complexity.LogProg.AReach, Complexity.LogProg.AReach.single, Complexity.LogProg.AReach.trans, Complexity.LogProg.astep}
\difficulty{1}
Let $A$ be a program, let an oracle and input be fixed, and let $G$ be a predicate on
configurations. If $G$ holds at a configuration $a$ and there is a finite run from the
one-step successor of $a$ to a configuration $c$ whose configurations before the
endpoint all satisfy $G$, then there is a finite run from $a$ to $c$ whose
configurations before the endpoint all satisfy $G$.
\end{lemma}

\begin{lemma}[Prepending one step to guarded halting]
\lean{Complexity.LogProg.AHalt.step}
\label{Complexity.LogProg.AHalt.step}
\leanok
\uses{Complexity.LogProg.AConf, Complexity.LogProg.AHalt, Complexity.LogProg.AReach.halt, Complexity.LogProg.AReach.single, Complexity.LogProg.astep}
\difficulty{1}
Let $A$ be a register program with oracle, let $x$ be an input, and let $G$ be a
predicate on configurations. If $G$ holds of a configuration $a$ and the run starting
from the one-step successor of $a$ halts with answer $r$ while satisfying $G$ at every
configuration before halting, then the run starting from $a$ also halts with answer $r$
and satisfies $G$ at every configuration before halting.
\end{lemma}

\begin{lemma}[Halting after a return instruction]
\lean{Complexity.LogProg.AHalt.ret}
\label{Complexity.LogProg.AHalt.ret}
\leanok
\uses{Complexity.LogProg.AHalt, Complexity.LogProg.arun, Complexity.LogProg.astep}
\difficulty{4}
Let $A$ be a register program with oracle, let $x$ be an input, and let $G$ be a
predicate on configurations. If the instruction of $A$ at state $l$ is to return the
Boolean value $r$, and $G$ holds of the configuration with state $l$, register values
$v$, and no recorded result, then the run from that configuration halts with answer $r$
and satisfies $G$ at every configuration before halting.
\end{lemma}

\begin{lemma}[Monotonicity of guarded reachability]
\lean{Complexity.LogProg.AReach.mono}
\label{Complexity.LogProg.AReach.mono}
\leanok
\uses{Complexity.LogProg.AConf, Complexity.LogProg.AReach}
\difficulty{3}
Let $A$ be a register program with oracle, let $x$ be an input, and let $G$ and $G'$ be
predicates on configurations. If a run reaches configuration $b$ from configuration $a$
with $G$ holding at every configuration before reaching $b$, and every configuration
satisfying $G$ also satisfies $G'$, then it reaches $b$ from $a$ with $G'$ holding at
every configuration before reaching $b$.
\end{lemma}

\begin{lemma}[Monotonicity of guarded halting]
\lean{Complexity.LogProg.AHalt.mono}
\label{Complexity.LogProg.AHalt.mono}
\leanok
\uses{Complexity.LogProg.AConf, Complexity.LogProg.AHalt}
\difficulty{3}
Let $A$ be a register program with oracle, let $x$ be an input, let $G$ and $G'$ be
predicates on configurations, and let $a$ be an initial configuration. If the run from
$a$ halts with answer $r$ and every configuration before halting satisfies $G$, and
every configuration satisfying $G$ also satisfies $G'$, then the run from $a$ halts with
answer $r$ and every configuration before halting satisfies $G'.$
\end{lemma}

\begin{lemma}[Monotonicity of Clean Runs]
\lean{Complexity.LogProg.cleanRun_mono}
\label{Complexity.LogProg.cleanRun_mono}
\leanok
\uses{Complexity.LogProg.CleanRun, Turing.MultiTapeTM}
\difficulty{3}
Let $D$ be a multi-tape Turing machine with Boolean work-tape symbols, initial state
$q$, input word $V$, and Boolean output $b$. Suppose there is a time $T$ at which the
machine has halted, output exactly the one-symbol word $[b]$, and has blank contents on
every work tape with every work-tape head at the origin; suppose also that every
work-tape head, at every time $t\le T$, has absolute position at most $B$. If $B\le B'$,
then the same machine has a clean run with bound $B'$.
\end{lemma}

\begin{definition}[Program precondition]
\lean{Complexity.LogProg.PreS}
\label{Complexity.LogProg.PreS}
\leanok
\uses{Complexity.LogProg.AConf, Complexity.LogProg.ARM, Complexity.LogProg.ValidPair, Complexity.LogProg.ValidPlain}
Let $A$ be a program with $m$ registers, and let $x$ be an input bit string. A
configuration satisfies the precondition when, if its control state is an instruction
comparing two registers, those registers are distinct; if it compares a register with
the input payload, the input has the prescribed valid plain encoding; if it compares a
register with either component of a paired payload, the input has the prescribed valid
paired encoding; and if it is a decider call, the list of argument registers has no
repetitions. All other configurations satisfy the precondition.
\end{definition}

\begin{lemma}[Upper bound for virtual word length]
\lean{Complexity.LogProg.length_vword_le}
\label{Complexity.LogProg.length_vword_le}
\leanok
\uses{Complexity.LogProg.Seg, Complexity.LogProg.length_render, Complexity.LogProg.rlen, Complexity.LogProg.vword}
\difficulty{4}
For any finite list of segments, the length of the virtual word formed by rendering them
in order and placing the separator $01$ between consecutive segments is at most the sum
of their rendered lengths plus twice the number of segments.
\end{lemma}

\begin{lemma}[First components of argument segments]
\lean{Complexity.LogProg.mem_argSegs}
\label{Complexity.LogProg.mem_argSegs}
\leanok
\uses{Complexity.LogProg.Seg, Complexity.LogProg.argSegs}
\difficulty{4}
Let $ws$ be a list of Boolean words, and let $sg$ be a segment belonging to the list of
argument segments formed from $ws$. Then the word component of $sg$ occurs in $ws$.
\end{lemma}

\begin{lemma}[Bound on Total Argument Segment Length]
\lean{Complexity.LogProg.sum_rlen_argSegs_le}
\label{Complexity.LogProg.sum_rlen_argSegs_le}
\leanok
\uses{Complexity.LogProg.argSegs, Complexity.LogProg.mem_argSegs, Complexity.LogProg.rlen}
\difficulty{4}
Let $ws$ be a finite list of bit strings, and suppose every string in $ws$ has length at
most $W$. Form argument segments by marking each word for doubling except the final
word, which is left undoubled. The sum of the segment lengths, counting doubled words
twice, is at most $|ws|\,(2W)$.
\end{lemma}

\begin{lemma}[Length bound for encoded call words]
\lean{Complexity.LogProg.length_callSegs_le}
\label{Complexity.LogProg.length_callSegs_le}
\leanok
\uses{Complexity.LogProg.CallSpec, Complexity.LogProg.Mode.seg0, Complexity.LogProg.callSegs, Complexity.LogProg.length_argSegs, Complexity.LogProg.length_vword_le, Complexity.LogProg.rlen, Complexity.LogProg.sum_rlen_argSegs_le, Complexity.LogProg.vword}
\difficulty{5}
Let $m,d$ be nonnegative integers, let $\Lambda$ be a type, and let a call specification
have at most $m$ registers, select a decider among $d$ choices, choose a mode and an
ordered list of registers, and specify successor states for positive and negative
answers. For an input bit string $x$, assign each register $r$ a bit string $V(r)$, and
suppose every assigned string has length at most $W$ and the selected register list
contains no repetitions. Then the virtual call word formed from the mode-dependent first
segment, the selected register strings (with the prescribed doubling), and separators
has length at most $2|x|+2+m(2W+2)$, where $|x|$ denotes the length of $x$.
\end{lemma}

\begin{theorem}[Logspace decidability by an oracle register program]
\lean{Complexity.LogProg.arm_decides}
\label{Complexity.LogProg.arm_decides}
\leanok
\uses{Complexity.LOGSPACE, Complexity.LogProg.AHalt, Complexity.LogProg.ARM, Complexity.LogProg.Pre, Complexity.LogProg.PreS, Complexity.LogProg.armProg, Complexity.LogProg.arm_space, Complexity.LogProg.arun, Complexity.LogProg.bankK, Complexity.LogProg.bankStart, Complexity.LogProg.bankTM, Complexity.LogProg.bank_cleanRun, Complexity.LogProg.callSegs, Complexity.LogProg.cleanRun_mono, Complexity.LogProg.compileFinTM, Complexity.LogProg.length_bits_mono, Complexity.LogProg.length_callSegs_le, Complexity.LogProg.log_poly_bound, Complexity.LogProg.vword, Complexity.logSpace, Turing.FinTM, Turing.FinTM.DecidesInSpace, Turing.MultiTapeTM.indicator}
\difficulty{7}
Let $L$ be a language of bit strings, let $A$ be a register program with $m$ registers
and finitely many states, and let $\ell_0$ be its initial state. Suppose there are $d$
oracle languages $A_0,\ldots,A_{d-1}$, each decidable in logarithmic space, and a
natural number $K$, such that for every input $x$, the run starting at $\ell_0$ with all
registers zero halts with the correct membership answer for $L$. At every configuration
before halting, the program's side conditions hold and, for each register $r$, the
binary representation of its value plus one has length at most $K(\lfloor\log_2
|x|\rfloor+1)$. Then $L$ is decidable in logarithmic space.
\end{theorem}

%%% added by the blueprint pipeline
\subsection*{Space Complexity --- Machines --- ARMProof: further results}

\begin{definition}[Instruction relabeled into a larger machine]
\lean{Complexity.LogProg.Ins.embed}
\label{Complexity.LogProg.Ins.embed}
\leanok
\uses{Complexity.LogProg.Ins}
Let $\rho$ map the $m$ registers of a register machine into $m'$ registers, $\sigma$ map its
$d$ oracle deciders into $d'$ deciders, $\iota$ map its labels $\Lambda$ into labels
$\Lambda'$, let $\kappa:\{\mathrm{true},\mathrm{false}\}\to\Lambda'$ assign an exit label to
each answer, and let $r_0$ be one of the $m'$ registers. The embedding of an instruction
renames every register it mentions by $\rho$ (elementwise on a call's argument list), its
decider by $\sigma$ and every target label by $\iota$, keeping the kind of instruction, except
that the instruction ``halt with answer $b$'' becomes a zero test of $r_0$ whose two branches
both jump to $\kappa(b)$.
\end{definition}

\begin{definition}[Register extension along a register map]
\lean{Complexity.LogProg.RegExt}
\label{Complexity.LogProg.RegExt}
\leanok
Let $\rho$ map $m$ registers into $m'$ registers, let $e$, $v'$ assign naturals to the $m'$
registers and $v$ to the $m$ registers. Then $v'$ extends $v$ along $\rho$ with background $e$
when $v'(\rho(r))=v(r)$ for every register $r$, and $v'(r')=e(r')$ for every register $r'$ not
in the range of $\rho$.
\end{definition}

\begin{lemma}[Register writes preserve extensions]
\lean{Complexity.LogProg.RegExt.update}
\label{Complexity.LogProg.RegExt.update}
\leanok
\uses{Complexity.LogProg.RegExt}
\difficulty{4}
Let $\rho$ be an injective map from $m$ registers into $m'$ registers, let $e,v'$ assign
naturals to the $m'$ registers and $v$ to the $m$ registers, and suppose $v'$ extends $v$ along
$\rho$ with background $e$ (that is, $v'(\rho(r))=v(r)$ for all $r$ and $v'$ agrees with $e$
off the range of $\rho$). Then for every register $r$ and every $n\in\mathbb N$, the assignment
$v'$ with $\rho(r)$ set to $n$ extends $v$ with $r$ set to $n$, along $\rho$ with background
$e$.
\end{lemma}

\begin{lemma}[Syntactic preconditions from the instruction]
\lean{Complexity.LogProg.preS_of_ins}
\label{Complexity.LogProg.preS_of_ins}
\leanok
\uses{Complexity.LogProg.ARM, Complexity.LogProg.PreS, Complexity.LogProg.ValidPair, Complexity.LogProg.ValidPlain}
\difficulty{3}
Let $A$ be a register machine with $m$ registers, $d$ oracle deciders and labels $\Lambda$, let
$x$ be a Boolean input, $\ell$ a label, $v$ a register assignment and $a$ an optional answer.
Suppose that if the instruction at $\ell$ is an equality test of registers $r,s$ then
$r\ne s$; if it is an oracle call then its argument registers are pairwise distinct; if it
compares a register with the payload of a plain input then $x$ has the form
$\langle 1^n,w\rangle$ with $w$ not ending in $0$; and if it compares a register with
either payload of a pair input then $x$ has the form $\langle 1^n,\langle u,w\rangle\rangle$
with neither $u$ nor $w$ ending in $0$. Then the configuration $(\ell,v,a)$ satisfies the syntactic
step preconditions of $A$ on $x$.
\end{lemma}

\begin{lemma}[Preconditions carry over to an embedded block]
\lean{Complexity.LogProg.preS_embed}
\label{Complexity.LogProg.preS_embed}
\leanok
\uses{Complexity.LogProg.Ins.embed, Complexity.LogProg.PreS}
\difficulty{3}
Let $A$ be a register machine with $m$ registers, $d$ deciders and labels $\Lambda$, and $D$ one
with $m'$ registers, $d'$ deciders and labels $\Lambda'$. Let $\rho$ be an injective map of
registers, $\sigma$ a map of deciders, $\iota$ a map of labels, $\kappa$ an exit label for each
answer and $r_0$ a register of $D$, and suppose that for every label $\ell$ the instruction of
$D$ at $\iota(\ell)$ is the instruction of $A$ at $\ell$ with registers renamed by $\rho$,
deciders by $\sigma$, labels by $\iota$, and an answer $b$ replaced by a zero test of $r_0$
jumping to $\kappa(b)$ on both branches. If the live configuration of $A$ at $\ell$ with
registers $v$ satisfies $A$'s syntactic step preconditions on a Boolean input $x$, then for
every register assignment $v'$ of $D$ the live configuration of $D$ at $\iota(\ell)$ with
registers $v'$ satisfies $D$'s syntactic step preconditions on $x$.
\end{lemma}

\begin{lemma}[One step of an embedded block]
\lean{Complexity.LogProg.astep_embed}
\label{Complexity.LogProg.astep_embed}
\leanok
\uses{Complexity.LogProg.CallSpec, Complexity.LogProg.Ins.embed, Complexity.LogProg.RegExt, Complexity.LogProg.RegExt.update, Complexity.LogProg.astep, Complexity.LogProg.astep_call, Complexity.LogProg.astep_clr, Complexity.LogProg.astep_dec, Complexity.LogProg.astep_half, Complexity.LogProg.astep_inc, Complexity.LogProg.astep_jeq, Complexity.LogProg.astep_jeqFst, Complexity.LogProg.astep_jeqIn, Complexity.LogProg.astep_jeqSnd, Complexity.LogProg.astep_jodd, Complexity.LogProg.astep_jz, Complexity.LogProg.astep_ret, Complexity.LogProg.callSegs, Complexity.LogProg.callSegs_congr}
\difficulty{5}
Let $A$ be a register machine with $m$ registers, $d$ deciders and labels $\Lambda$, and $D$ one
with $m'$ registers, $d'$ deciders and labels $\Lambda'$; let $\rho$ be an injective map of
registers, $\sigma$ a map of deciders, $\iota$ a map of labels, $\kappa$ an exit label for each
answer and $r_0$ a register of $D$, such that for every label $\ell$ the instruction of $D$ at
$\iota(\ell)$ is the instruction of $A$ at $\ell$ with registers renamed by $\rho$, deciders by
$\sigma$, labels by $\iota$, and an answer $b$ replaced by a zero test of $r_0$ jumping to
$\kappa(b)$ on both branches. Let $\mathcal O$ be an oracle for the $d'$ deciders, $x$ a
Boolean input, $\ell$ a label at which $A$ has neither of the two input-shape validation
instructions, and let $v'$ extend $v$ along $\rho$ with background $e$ ($v'(\rho(r))=v(r)$
for all $r$, $v'=e$ off the range of $\rho$). Then either one step of $A$ on $x$ with oracle
$\mathcal O\circ\sigma$ from the live configuration $(\ell,v)$ leads to a live configuration
$(\ell_2,v_2)$ and one step of $D$ on $x$ with oracle $\mathcal O$ from $(\iota(\ell),v')$
leads to a live configuration $(\iota(\ell_2),v_2')$ where $v_2'$ extends $v_2$ along $\rho$
with background $e$; or, for some answer $b$, that step of $A$ halts with answer $b$ and
registers $v$, and that step of $D$ leads to the live configuration $(\kappa(b),v')$.
\end{lemma}

\begin{theorem}[A halting run of an embedded block]
\lean{Complexity.LogProg.AHalt.embed}
\label{Complexity.LogProg.AHalt.embed}
\leanok
\uses{Complexity.LogProg.AConf, Complexity.LogProg.AHalt, Complexity.LogProg.AReach, Complexity.LogProg.AReach.refl, Complexity.LogProg.AReach.step, Complexity.LogProg.Ins.embed, Complexity.LogProg.RegExt, Complexity.LogProg.arun, Complexity.LogProg.arun_halted, Complexity.LogProg.arun_succ, Complexity.LogProg.arun_zero, Complexity.LogProg.astep_embed}
\difficulty{5}
Let $A$ be a register machine with $m$ registers, $d$ deciders and labels $\Lambda$, and $D$ one
with $m'$ registers, $d'$ deciders and labels $\Lambda'$; let $\rho$ be an injective map of
registers, $\sigma$ a map of deciders, $\iota$ a map of labels, $\kappa$ an exit label for each
answer and $r_0$ a register of $D$, such that for every label $\ell$ the instruction of $D$ at
$\iota(\ell)$ is the instruction of $A$ at $\ell$ with registers renamed by $\rho$, deciders by
$\sigma$, labels by $\iota$, and an answer $b$ replaced by a zero test of $r_0$ jumping to
$\kappa(b)$ on both branches; and suppose $A$ contains no input-shape validation instruction.
Let $\mathcal O$ be an oracle for the $d'$ deciders, $x$ a Boolean input, $G$ and $G'$
predicates on configurations of $A$ and of $D$, and $e$ a background assignment for $D$'s
registers, such that whenever a live configuration $(\ell,v)$ of $A$ satisfies $G$ and $v'$
extends $v$ along $\rho$ with background $e$, the live configuration $(\iota(\ell),v')$ of $D$
satisfies $G'$. If $A$ on $x$ with oracle $\mathcal O\circ\sigma$, started at the live
configuration $(\ell,v)$, halts with answer $b$ with every configuration before halting
satisfying $G$, and $v'$ extends $v$ along $\rho$ with background $e$, then there are
assignments $v_2$ and $v_2'$ such that $(\ell_2,v_2)$ satisfies $G$ for some label $\ell_2$,
$v_2'$ extends $v_2$ along $\rho$ with background $e$, and $D$ on $x$ with oracle $\mathcal O$
runs from $(\iota(\ell),v')$ to the live configuration $(\kappa(b),v_2')$ with every
configuration before the end satisfying $G'$.
\end{theorem}
